\documentclass[10pt,a4paper,twocolumn]{book}
%PREAMBULO
	
	\usepackage[utf8]{inputenc} % Acentuação 
	\usepackage{multirow} %Mesclagem vertical de linhas
	\usepackage{graphicx} %Inserção de figuras
	\usepackage{subfigure}	
	\usepackage[portuguese]{babel} %Tradução dos labels
	\usepackage{amsfonts} %Símbolos matemáticos
		\usepackage{amsmath} %Matrizes e outros símbolos		
		\usepackage{amsthm} %Ambientes para detalhes de teoremas
	\usepackage[ruled,linesnumbered,portuguese,onelanguage]{algorithm2e} %Algoritmos pseudocódigo
	\usepackage{float} %flutuaçao de elementos/ posição relativa
	\usepackage{listings} %Inserção de códigos 
	\usepackage{tikz,pgfplots} %Pacote gráfico
		% Tikz importado
		\pgfplotsset{compat=1.15}
		\usepackage{mathrsfs}
		\usetikzlibrary{arrows}
	\usepackage{natbib}	 %Referências bibliográficas
	\usepackage[top=2cm, bottom=2cm, left=2.5cm, right=2.5cm]{geometry} %Dimensionar estilos da página
		
	\newtheorem{teorema}{Teorema}
	\newtheorem{corolario}{Corolário}[teorema]
	\newtheorem{lema}{Lema}[corolario]
	
\begin{document} 
%CORPO DO DOCUMENTO
	
	\bibliographystyle{unsrtnat} %Padrão da bibliografia/Estilo
	\tableofcontents
	
	\chapter{Minicurso de \LaTeX}
	
	Hello World! \\ 
	Introdução!
	Problemas na acentua\c{c}\~ao. Assim que se resolve no modo hard!
	
	\section{Fontes}
		\begin{enumerate}
			\item Tamanho
				\begin{itemize}
					\item {\tiny Pequena} 
					\item {\small Texto } 
					\item {\Large Oláaaa}  
					\item {\LARGE Big }
					\item {\Huge Oahdsad}
				\end{itemize}
			\item Estilo
				\begin{itemize}
					\item \textbf{Negrito} ou {\bf negrito mais negrito}
					\item \textit{Itálico} ou {\it Itálico mais mais}
					\item \underline{Sublinhado}
				\end{itemize}
			
		\end{enumerate}
	
	
	
	\section{Titulo da seção}
		Esta é minha primeira seção.
		\subsection{Subseção}
			Uma coleção de itens:
			\begin{itemize}
				\item Item 1
				\item Item 2
			\end{itemize}
			Quero uma enumeração
		
			\begin{enumerate}
				\item Bom
				\begin{enumerate}
					\item Um subitem
					\begin{itemize}
						\item Um subsubitem
					\end{itemize}
				\end{enumerate}
				\item Olá
				\item[labelPerson] Último item 
			\end{enumerate}
			
			\subsubsection{Subsubseção}
 	\section{Segunda Seção}
		Vamos trocar uma seção pela outra?
	
	\section{Outros elementos}
   		\begin{table}[h]
			\caption{Minha primeira tabela}
			\centering
			\begin{tabular}{ c | l | r } \hline
				\multicolumn{3}{c}{Texto mesclado} \\ \hline
				Centralizado &esquerda &direita \\ \hline
				linha 2 & & \\ 
				linha3 & & \\ \hline
			\end{tabular}
		\end{table}
		
		\begin{table}[h]
			\caption{Tabela com mesclagem de linhas}
			\centering
			\begin{tabular}{ c | l | r } \hline
				\multicolumn{3}{c}{Texto mesclado} \\ \hline
				Centralizado &esquerda &direita \\ \hline
				linha 2 & & \\ 
				linha3 & & \\ \hline
				\multirow{2}{*}{Celulas mescladas} & & \\ \hline
				& & \\
				& & \\ \hline
				linha 6 & &\\ \cline{2-3}
				linha 7 & &\\ \hline
			\end{tabular}
		\end{table}
		
	\section*{Vamos inserir uma imagem?}
	
		\begin{figure}[h]
			\caption{Legenda huehue}
			\includegraphics[scale=0.2]{figuras/encomp-logo.png}
		\end{figure}
	
		\begin{figure}[h]
			\caption{Subfiguras}
			\label{fig:subfiguras}
			\subfigure[Figura a]{\includegraphics[scale=0.1]{figuras/encomp-logo.png} \label{subfig:logo}}
			\subfigure[Figura  b ]{\includegraphics[scale=0.18]{figuras/joao-crazy.jpg}}
		\end{figure}
	
		A figura \ref{fig:subfiguras} mostra duas fig.
%	\part{Uma parte}
	
		\section{MathLovers} 
			Uma seção para morrer de amores.
			Que a $f = m \cdot a $ esteja com vocês.
			Na forma vetorial, fica $$\vec{f} = m \cdot \vec{a}.$$
			
			Num triângulo qualquer, vale a Lei dos Cossenos
			\begin{equation}
				a^2 + b^2 - 2\cos(\theta) = c^2\mbox{, com } \theta = \hat{A}.
				\label{eq:leiCos}
			\end{equation}
			Em particular, num triângulo retângulo, $\theta = \frac{\pi}{2}$.
			Daí
			\begin{math}
				 a^2 + b^2 = c^2
			\end{math}
			
			Algumas identidades trigonométricas...
			
			\begin{eqnarray}
				\sin^2(x) + \cos^2(x) &= &1 \label{eq: relTrigFund}\\ 
				1 + \tan^2(x) &= &\sec^2(x) \\ \nonumber
				\cos(2\theta) &= &\cos^2(\theta) - \sin^2(\theta)
			\end{eqnarray}
		
			Função contínua
			\begin{math}
				\lim\limits_{x \rightarrow a} f(x) = f(a)
			\end{math}
			
			Alguns dos conjuntos numéricos mais estudados
			$$ \mathbb{N} \subset \mathbb{Z} \subset \mathbb{R}$$
			Exemplo: $ \mathbb{R}_+ = \{ x \in \mathbb{R} | x \ge 0\}$
			
			Uma transformação linear $T$ importantíssima na computação é a rotação, definida como  \\$T:\mathbb{R}^2 \longrightarrow \mathbb{R}^2$ cuja matriz associada é 
			
			\begin{equation}
				\begin{bmatrix}
					cos(\alpha) &\sin(\alpha) \\
					-\sin(\alpha) &\cos(\alpha)
				\end{bmatrix}
			\end{equation}
			
			É fácil ver que 
			
			\begin{equation}
				\begin{vmatrix}
					cos(\alpha) &\sin(\alpha) \\
					-\sin(\alpha) &\cos(\alpha)
				\end{vmatrix}
				= 1, \text{ por \ref{eq: relTrigFund}. }
			\end{equation}
			
			Outro exemplo é 
			\begin{equation}
				|x| = \sqrt{x^2} = 
					\begin{cases}
						x, &\text{se } x \ge 0 \\
						-x, &\text{se } x < 0 
					\end{cases}
			\end{equation}
			
			\begin{equation}
				(x+y)^n = \sum_{i=0}^{n} 
					\begin{pmatrix}
						n \\
						i
					\end{pmatrix}
				x^i\cdot y^{n-i}
			\end{equation}
			
	\subsection{Para brincar}
		\begin{teorema}
			Aqui vai um teorema.
		\end{teorema}
	
		\begin{teorema}[Teorema Fundamental do Cálculo]
			Seja $f$ uma função ...
			Então 
			\begin{equation}
			\frac{d}{dx}\int_{a}^{x}f(t)dt = f(x)
			\end{equation}
		\end{teorema}
	
		\begin{corolario}
			Ui! Olá sou um Corolário
		\end{corolario}
			
		\begin{lema}
			OOps
		\end{lema}
		
		\begin{teorema}[Teorema de Chicó] Tudo é porque é porque é mesmo.
			\begin{proof}
				De fato, é porque é, porque é mesmo.
			\end{proof}
		\end{teorema}
	Conforme demonstrado, este resultado é válido para tudo
	
	
	\section{Computarias}
		Vamos inserir um algoritmo:
		
		\begin{algorithm}[h]
			\caption{{\sc Merge }}
			\DontPrintSemicolon % Some LaTeX compilers require you to use \dontprintsemicolon instead
			\KwIn{A[], p, q, r : inteiro}
			\KwOut{Sequência A[p, $\ldots$ , r-1] ordenada}
			$i, j, k,$ temp[$1, \ldots, r-p$]: inteiro 
			
			$i \gets p$;\\
			$j \gets q$;\\
			$k \gets 1$;\\
			
			\While{$i < q$ \textbf{e} $j < r$ }
			{
				\label{alg:laco}
				\If{A[$i$] $\leq$ A[$j$] }{temp[$k$] $\gets$ A[$i$];\\ $i \gets i + 1$; }
				\Else{temp[$k$] $\gets$ A[$j$];\\ $j \gets j + 1$; }
				$k \gets k + 1$;
			}
			
			\While{$i < q$ }
			{
				temp[$k$] $\gets$ A[$i$];\;
				$i \gets i + 1$;\;
				$k \gets k + 1$;
			}
			
			\While{$j < r$ }
			{
				temp[$k$] $\gets$ A[$j$];\; 
				$j \gets j + 1$;\;
				$k \gets k + 1$;
			}			
			
			A[$p , \ldots, r-1$] $\gets$ temp[$1, \ldots, r-p$]\;
			\Return{A}\;
			\label{algo:selectionsort}
		\end{algorithm}	
		
		No laço da linha \ref{alg:laco}...
		Ahh mas quero inserir meu código em outra linguagem...
		
\lstset{language=C}			
Aqui está um pequeno código que calcula e exibe o fatorial de um número natural em C

\begin{lstlisting}
#include <stdio.h>

#define pi 3.1415926535

unsigned int fatorial( unsigned int x )
{
 return ( x > 1 ) ? x*fatorial( x-1 ) : 1;
}

int main()
{
 int pi_f = (int)pi; 
 printf("%u \n", \
 fatorial( (unsigned int)pi_f )  );
 return 0;
}
\end{lstlisting}

Agora um código em MATLAB inserindo de arquivo externa.

\lstinputlisting[language=MATLAB, breaklines=true]{codigos/pagerank.m}

\section{Organização externa}
	Agora estamos trabalhando com modularização...



	
	
	\bibliography{bibliografia}
		\nocite{livro:joao}
		\nocite{conf:alckmin}
\end{document}